% 20260703_Assingment_Probability_Statistics_A
\documentclass[dvipdfmx]{jsreport}
%preamble
\usepackage{soupreamble}

\title{確率及び統計学特論A 第5回レポート課題}
\author{M1 6012 川村聡一郎}
\date{}


\begin{document}
\maketitle

\begin{tcolorbox}
	$\mu$の分布関数を$F$とし，$F$の不連続点を$A_\mu$とする．このとき$A_\mu$は高々加算であることを示せ．
%Hint: 分布関数をFとして，$A_{\mu, n}= \{x\in \R; F(x)- F(x-)> 1/n\}$とおくと
% #(A_{\mu, n}の元)\le n(背理法)
\end{tcolorbox}

\begin{souanswer} %答案はここに書く
	分布関数$F$は単調非減少であり，$0\le F(x)\le 1$である．よって任意の点$x$で左極限が存在し，$F(x)- F(x- 0)\ge 0$である．\\
	各$n\in \N$に対して，$A_{\mu, n}= \{x\in \R; F(x)- F(x^-)> 1/n\}$とおく．\\
	$F(\infty)- F(-\infty)\ge 1$であることから幅$1/n$を超える点は高々$n$個しか存在せず，$\#(A_{\mu, n})\le n$をみたす．\\
	元の不連続点全体の集合$A_\mu$は$A_\mu= \bigcup_{n= 1}^\infty A_{\mu, n}$と表せ，有限集合の可算無限個の和集合は高々可算集合であるから，$A_\mu$も高々可算．
\end{souanswer}
\end{document}